\documentclass[10pt,a4paper]{amsart}
\usepackage[margin=19mm]{geometry}
\usepackage{amsmath,amssymb,mathtools}
\usepackage[scr=rsfs,cal=euler]{mathalfa}
\usepackage{xfrac}
\usepackage[unicode,hidelinks,bookmarks=false]{hyperref}
\newcommand{\ie}{i.e.\ }
\newcommand{\cf}{cf.\ }
\newcommand{\resp}{respectively\ }
\newcommand{\Subsection}{\subsection}
\providecommand{\nobreakdash}{\nobreak\hbox{-}}
\setlength{\emergencystretch}{2em}
\multlinegap0pt
\usepackage{amssymb}
\usepackage[shortlabels]{enumitem}
\newtheorem{prop}{Proposition}
\newtheorem{theorem}{Theorem}
\usepackage{cleveref}
\crefname{prop}{Proposition}{Propositions}
\crefname{theorem}{Theorem}{Theorems}
\newcommand{\M}{{\mathcal M}}
\newcommand{\R}{{\mathbb R}}
\newcommand{\Z}{{\mathbb Z}}
\newcommand{\la}{\langle}
\newcommand{\p}{\partial}
\newcommand{\ra}{\rangle}
\newcommand{\we}{\wedge}
\title{Approximate controllability on the group of volume-preserving diffeomorphisms}
\author{Andrei Agrachev, Bettina Kazandjian}
\date{Source: arXiv:2601.16939v1 (CC BY 4.0)}
\begin{document}
\maketitle

\noindent\emph{Source and licence.} Approximate controllability on the group of volume-preserving diffeomorphisms. Version arXiv:2601.16939v1 (CC BY 4.0), distributed by arXiv under the Creative Commons Attribution 4.0 International licence, http://creativecommons.org/licenses/by/4.0/, as recorded in the arXiv licence metadata for this exact version and shown on the abstract page https://arxiv.org/abs/2601.16939v1. Full licence notice: \texttt{licenses/arxiv-2601.16939-CC-BY-4.0.txt}.\par\smallskip

\noindent\textit{Editorial scope.} Counted unit: the article's theorem span3 (source0.tex lines 706--712). The article's complete original proof of it is delivered with it (source0.tex lines 887--1000, one \texttt{proof} environment of 114 lines, which the article places away from the statement). No mathematics is written, repaired or re-derived: the statement and the proof are the article's own lines, in the article's own notation, and no retained line is substituted or re-flowed.  Retained uncounted as the source-local context the proof needs: lines 639--654; lines 657--662; lines 716--737; lines 791--798; lines 849--858.  Nothing was omitted from any retained span, so the package carries no omission record.  No retained line contains a citation, so the wrapper carries no bibliography and no bibliography entry.  No retained line contains a figure environment, an image-inclusion command or drawing code, so no image is delivered and the package declares no asset dependency.  Editorial formatting only, in addition: an AMS \texttt{amsart} preamble that restates the article's own declarations and macros used by the retained lines, namely the environments \texttt{prop}, \texttt{theorem} on separate counters, together with the standard packages those lines need.  The retained environments print in this document as Proposition 1, Theorem 1 in order of appearance, whereas the article prints the same items with the numbering of its own full document.  The complete native source of the article as distributed in the arXiv e-print for this exact version is bundled unchanged as \texttt{sources/193279/source0.tex}.
\begin{prop} \label{formulas T3}
\begin{align*}
    1)\quad &[p_m\sin\la m,\cdot \ra,p_n\cos \la n,\cdot \ra]=\la m,p_n\ra p_m\cos \la m,\cdot \ra \cos \la n,\cdot\ra +\la n,p_m\ra p_n\sin \la m,\cdot \ra \sin \la n,\cdot \ra,  \\
    &[p_m\cos \la m,\cdot \ra,p_n \sin \la m,\cdot \ra]=-\la m,p_n\ra p_m\sin \la m,\cdot \ra \sin \la n,\cdot\ra -\la n,p_m\ra p_n\cos \la m,\cdot \ra \cos \la n,\cdot \ra, \\
    &[p_m\cos \la m,\cdot \ra,p_n \cos \la m,\cdot \ra]=-\la m,p_n\ra p_m\sin \la m,\cdot \ra \cos \la n,\cdot\ra +\la n,p_m\ra p_n\cos \la m,\cdot \ra \sin \la n,\cdot \ra, \\
    &[p_m\sin \la m,\cdot \ra,p_n \sin \la m,\cdot \ra]=-\la m,p_n\ra p_m\cos \la m,\cdot \ra \sin \la n,\cdot\ra -\la n,p_m\ra p_n\sin \la m,\cdot \ra \cos \la n,\cdot \ra. \\ \\
    2)\quad &[p_m\sin \la m,\cdot \ra,p_n\sin\la n,\cdot\ra ]-[p_m\cos \la m,\cdot \ra,p_n\cos \la n,\cdot \ra ]
    =(\la m,p_n\ra p_m - \la n,p_m\ra p_n)\sin \la m+n,\cdot \ra,\\
    &[p_m\cos \la m,\cdot \ra,p_n\sin\la n,\cdot\ra ]+[p_m\sin \la m,\cdot \ra,p_n\cos \la n,\cdot \ra ]
    =(\la m,p_n\ra p_m - \la n,p_m\ra p_n)\cos \la m+n,\cdot \ra.\\ \\
    3) \quad &[a_m\cos\la m,\cdot \ra+b_m\sin \la m,\cdot \ra,c_n\cos \la n,\cdot \ra+d_n \sin \la n,\cdot \ra]\\
        -&[-a_m\sin\la m,\cdot \ra+b_m\cos \la m,\cdot \ra,-c_n\sin \la n,\cdot \ra+d_n \cos \la n,\cdot \ra]\\
        &=(\la m,c_n\ra b_m+\la m,d_n\ra a_m- \la n,b_m \ra c_n - \la n,a_m \ra d_n)\cos \la m+n,\cdot \ra \\
        +&(\la m,d_n\ra b_m-\la m,c_n \ra a_m -\la n,b_m \ra d_n+\la n,a_m \ra c_n)\sin \la m+n,\cdot \ra.
    \end{align*}
\end{prop}

\begin{theorem} \label{theorem : appendix T3 span1}
     Let $\# \M_f < \infty$. If $\mathrm{span}\M_f$ is of dimension 1 or if ($\mathrm{span}\M_f$ is of dimension 2 and $\mathrm{span}\left\{a_m,b_m \mid m \in \M_f\right\}$ is of dimension 1), then 
        \begin{equation*}
            \mathrm{Lie}\left\{f+u \mid u \in \R^3\right\}=\mathrm{span}\left\{(\alpha a_m + \beta b_m)\cos \la m,\cdot \ra + (-\beta a_m + \alpha b_m)\sin \la m,\cdot \ra,\p_x,\p_y,\p_z \mid m\in \M_f, \alpha,\beta \in \R\right\},
        \end{equation*}
\end{theorem}

\begin{prop} \label{adding modes : different angles}
    Let $m\in \M_f$ be such that $m\neq 0$. If there exists $a_m,b_m,a_m',b_m'$ constant vector fields such that $a_m,a_m'\neq 0$,
    \begin{equation*}
        \left\{\begin{array}{cc}
           a_m\cos \la m,\cdot\ra+b_m \sin \la m,\cdot \ra \in \mathfrak{L}_f,    \\
        a_m'\cos\la m,\cdot \ra+b_m'\sin \la m,\cdot \ra\in \mathfrak{L}_f,
        \end{array}\right.
    \end{equation*}
    and such that $\begin{pmatrix}
            a'_m \\
            b'_m
        \end{pmatrix}\neq \begin{pmatrix}
            \alpha & \beta\\
            -\beta & \alpha 
        \end{pmatrix}\begin{pmatrix}
            a_m\\
            b_m
        \end{pmatrix}$ for every $\alpha,\beta\in \R$, then for every $p_m,q_m\in m^{\perp}$,
    \begin{equation*}
        p_m\cos \la m,\cdot \ra + q_m \sin \la m,\cdot \ra \in \mathfrak{L}_f.
    \end{equation*}
\end{prop}

\begin{prop} \label{generation of modes : from 2 non colinear modes}
    Lets $m,n\in \M_f$ be such that $m\we n \neq 0$. If there exist $a_m,a_m'\in m^{\perp}$ and $c_n,d_n \in n^{\perp}$ such that 
    \begin{align*}
        a_m\cos \la m,\cdot \ra \in \mathfrak{L}_f, \quad a_m'\cos \la m,\cdot \ra \in \mathfrak{L}_f, \quad
        c_n \cos \la n,\cdot \ra +d_n \sin \la n,\cdot \ra \in \mathfrak{L}_f,
    \end{align*}
    and such that $\mathrm{span}\left\{a_m,a_m',c_n\right\}=\R^3$, then $p_k\cos\la k,\cdot \ra + q_k\sin \la k,\cdot \ra \in \mathfrak{L}_f$ for every $p_k,q_k \in k^{\perp}$ and $k \in \la m,n\ra$, where $\la m,n\ra$ denotes the subgroup of $\Z^3$ generated by $m,n$.
\end{prop}

\begin{prop} \label{generation of every modes dimV0=3 dim I0=3}
    If $\mathrm{span}\left\{a_m,b_m \mid m \in \M_f \right\}=\R^3$ and $\mathrm{span}\M_f = \R^3$, and if there exists a non-zero mode $\ell \in \Gamma$, $p_{\ell},p_\ell'\in \ell^{\perp}$ such that $p_\ell\we p_\ell'\neq 0$ and such that 
    \begin{equation*}
        \left\{\begin{array}{ll}
             p_\ell \cos \la \ell,\cdot \ra \in \mathfrak{L}_f,\\
             p_\ell'\cos \la \ell,\cdot \ra \in \mathfrak{L}_f,
        \end{array}\right.
    \end{equation*}
    then $p_k\cos \la k,\cdot \ra +q_k \sin \la k,\cdot \ra \in \mathfrak{L}_f$, for every $p_k,q_k\in k^{\perp}$ and $k \in \Gamma$.
\end{prop}

\begin{theorem} \label{theorem : algebra dim3 span3}
    Let $\# \M_f < \infty$. If $\mathrm{span}\M_f=\R^3$ and $\mathrm{span}\left\{a_m,b_m \mid m \in \M_f\right\}=\R^3$,
    then  
    \begin{equation*}
        \mathrm{Lie}\left\{f+u \mid u \in \R^3\right\}=\mathrm{span} \left\{p_m\cos\la m,\cdot \ra,q_m\sin\la m,\cdot\ra,\p_x,\p_y,\p_z \mid m\in \Gamma,p_m,q_m \in m^{\perp}\right\}.
    \end{equation*}
\end{theorem}

\begin{proof}[Proof of \cref{theorem : algebra dim3 span3}.]
    If $\mathrm{span}\left\{a_m,b_m \mid m \in \M_f\right\}=\R^3$ and $\mathrm{span}\M_f=\R^3$, according to \cref{generation of every modes dimV0=3 dim I0=3}, it is sufficient to prove that there exists a non-zero mode $\ell \in \Gamma$ such that $p_\ell \cos \la \ell,\cdot \ra \in \mathfrak{L}_f $ for every $p_{\ell}\in \ell^{\perp}$. There are two possible situations: 
    \begin{enumerate}
        \item[1.] Either there exist $m,n\in \M_f$ such that $m\we n \neq 0$, 
        \begin{equation*}
            a_j\cos \la j,\cdot \ra +b_j \sin \la j,\cdot \ra \in \mathfrak{L}_f, \quad j\in \left\{m,n\right\},
        \end{equation*}
        and such that $\mathrm{span} \left\{a_m,b_m,a_n\right\}=\R^3$. 
        \item[2.] Or there exist $m,n,k\in \M_f$ such that $\mathrm{span} \left\{m,n,k\right\}=\R^3$, 
        \begin{equation}
            a_j \cos \la j,\cdot \ra +b_j \sin \la j,\cdot \ra \in \mathfrak{L}_f, \quad a_j \we b_j =0, \quad j \in \left\{m,n,k\right\},
        \end{equation}
        and such that $\mathrm{span} \left\{a_m,a_n,a_k\right\}=\R^3$.
    \end{enumerate}
    Let us consider both cases. \\
    \textbf{1.} Because $\mathrm{span} \left\{a_m,b_m,a_n\right\}=\R^3$, necessarily $a_m\we b_m \neq 0$. First we assume that $a_n \we b_n \neq 0$. According to \cref{theorem : appendix T3 span1}, for every $\alpha,\beta\in \R$,
    \begin{equation*}
        (\alpha a_j+\beta b_j)\cos \la j,\cdot \ra +(\alpha b_j-\beta a_j)\sin \la j ,\cdot \ra \in \mathfrak{L}_f \quad j \in \left\{m,n\right\}.
    \end{equation*}
    Thanks to an adapted choice of $\alpha,\beta$, we can assume that $b_m,b_n \in \mathrm{span} \left\{m\we n\right\}$. According to the third formula of \cref{formulas T3} applied with $(a_m,b_m,a_n,b_n)$, 
    \begin{equation*}
        p_{m+n}\cos \la m+n,\cdot \ra+q_{m+n}\sin \la m+n,\cdot \ra \in \mathfrak{L}_f, 
    \end{equation*}
    with 
    \begin{equation*}
        \left\{\begin{array}{ll}
             p_{m+n}=\la m,a_n\ra b_m  - \la n,a_m\ra b_n  \in \mathrm{span} \left\{m\we n\right\} \\
             q_{m+n}=- \la m,a_n \ra a_m + \la n,a_m \ra a_n. 
        \end{array}\right.
    \end{equation*}
    But it is also true that $a_n\cos \la -n,\cdot \ra-b_n \sin \la -n,\cdot\ra \in \mathfrak{L}_f$, so again we can apply the third formula of \cref{formulas T3} with $(p_{m+n},q_{m+n},a_{n},-b_n)$, and we obtain that $a_m'\cos \la m,\cdot \ra +b_m' \sin \la m,\cdot \ra \in \mathfrak{L}_f$, with 
    \begin{equation*}
    \left\{\begin{array}{ll}
         a_m'=-\la m,a_n \ra ^2 a_m  \\
         b_m'=-\la m,a_n \ra ^2 b_m + \la n,a_m \ra \la m,a_n \ra b_n.  
    \end{array}\right.
    \end{equation*}
    By assumption $a_n \we b_n \neq 0$, so $b_n \neq 0$. Moreover, $a_m,a_n\notin \mathrm{span} \left\{m \we n\right\}$, so $\la n,a_m \ra \la m,a_n \ra \neq 0$. Then $(a_m',b_m')\neq (\alpha a_m + \beta b_m, \alpha b_m -\beta a_m)$ for every $\alpha,\beta \in \R$ and $a_m,a_m' \neq 0$, so according to \cref{adding modes : different angles}, $p_m \cos \la m,\cdot \ra \in \mathfrak{L}_f$ for every $p_m \in m^{\perp}$. If we assume that $a_n \we b_n =0$, then $a_n \cos \la n,\cdot \ra \in \mathfrak{L}_f$. We can still assume that $b_m \we (m\we n)=0$, and because $\mathrm{span} \left\{a_m,b_m,a_n\right\}=\R^3$, necessarily $a_n \we b_m \neq 0$. Finally we obtain that
     \begin{equation*}
        p_{m+n}\cos \la m+n,\cdot \ra+q_{m+n}\sin \la m+n,\cdot \ra \in \mathfrak{L}_f, 
    \end{equation*}
    with 
    \begin{equation*}
        \left\{\begin{array}{ll}
             p_{m+n}=\la m,a_n\ra b_m   \\
             q_{m+n}=- \la m,a_n \ra a_m + \la n,a_m \ra a_n. 
        \end{array}\right.
    \end{equation*}
    But $m \we (m+n)\neq 0$, $a_m\we b_m \neq 0$, $p_{m+n}\we q_{m+n}\neq 0$ and $\mathrm{span} \left\{p_{m+n},q_{m+n},a_m\right\}=\R^3$, so according to the previous computations, $a_{m+n}\cos \la m+n,\cdot \ra \in \mathfrak{L}_f$ for every $a_{m+n}\in (m+n)^{\perp}$. 
    \\ \\
    \textbf{2.} If $a_j \cos \la j,\cdot \ra +b_j \sin \la j,\cdot \ra \in \mathfrak{L}_f$ and $a_j \we b_j=0$, then according to \cref{adding modes : different angles} $a_j \cos \la j,\cdot \ra \in \mathfrak{L}_f$. So in this case there exist $m,n,k\in \M_f$ such that $\mathrm{span} \left\{m,n,k\right\}=\R^3$ and such that
    \begin{equation*}
        a_m \cos \la m,\cdot \ra \in \mathfrak{L}_f, \quad a_n \cos \la n,\cdot \ra \in \mathfrak{L}_f, \quad a_k \cos \la k,\cdot \ra \in \mathfrak{L}_f,   
    \end{equation*}
    with $\mathrm{span} \left\{a_m,a_n,a_k\right\}=\R^3$. According to the second formulas of \cref{formulas T3}, $p_{m+n}\cos \la m+n,\cdot \ra\in \mathfrak{L}_f$ and $p_{n+k}\cos \la n+k,\cdot \ra \in \mathfrak{L}_f$ with 
    \begin{equation*}
        p_{m+n}= \la m,a_n \ra a_m - \la n,a_m \ra a_n, \quad p_{n+k}=\la n,a_k \ra a_n - \la k,a_n \ra a_k.
    \end{equation*}
    If $\la m,a_n\ra =0$ then $a_n \in \mathrm{span} \left\{ m\we n \right\}$ and necessarily $\la n,a_m\ra \neq 0$. So $p_{m+n}\neq 0$ and by symmetry $p_{n+k}\neq 0$.
    Then according to the same formulas, 
    \begin{equation*}
        p_{m,n,k}\cos \la m+n+k,\cdot \ra \in \mathfrak{L}_f, \quad p_{n,k,m}\cos \la n,k,m, \cdot \ra \in \mathfrak{L}_f,
    \end{equation*}
    with
    \begin{equation*}
        \left\{\begin{array}{ll}
             p_{m,n,k}=\la k,p_{m+n}\ra a_k -\la m+n,a_k \ra p_{m+n},  \\
             p_{n,k,m}=\la m,p_{n+k} \ra a_m - \la n+k,a_m\ra p_{n+k}.
        \end{array}\right.
    \end{equation*}
    If $\la m+n ,a_k \ra=0$ and $\la k,p_{m+n}\ra=0$, then $a_k,p_{m+n}\in \mathrm{span} \left\{(m+n)\we k\right\}$, and so $a_k \we p_{m+n}=0$. But $p_{m+n}\in \mathrm{span} \left\{a_m,a_n\right\}$ and $\mathrm{span} \left\{a_m,a_n,a_k\right\}=\R^3$, so necessarily $\la m+n,a_k \ra \neq 0$ or $\la k,p_{m+n}\ra \neq 0$ and $p_{m,n,k}\neq 0$. By symmetry $p_{n,k,m}\neq 0$. By computation we obtain that 
    \begin{equation*}
        \left\{\begin{array}{ll}
             p_{m,n,k}=-\la m+n,a_k\ra\la m,a_n \ra a_m + \la m+n,a_k\ra \la n,a_m \ra a_n + (\la m,a_n\ra \la k,a_m\ra -\la n,a_m \ra \la k,a_n\ra)a_k ,  \\ \\
             p_{n,k,m}=(\la m,a_n \ra \la n,a_k \ra-\la m,a_k \ra \la k,a_n \ra)a_m -\la n+k,a_m \ra \la n,a_k \ra a_n +\la n+k,a_m \ra \la k,a_n \ra a_k. 
        \end{array}\right.
    \end{equation*}
    If $p_{m,n,k}\we p_{n,k,m}=0$, then there exist a non-zero $\lambda\in \R$ such that 
    \begin{equation*}
        \left\{\begin{array}{ll}
             -\la m+n,a_k \ra \la m,a_n \ra = \lambda (\la m,a_n \ra \la n,a_k \ra - \la m,a_k \ra \la k,a_n \ra )  \\
              \la m+n,a_k \ra \la n,a_m \ra = -\lambda \la n+k,a_m \ra \la n,a_k \ra \\
              \la m,a_n \ra \la k,a_m \ra - \la n,a_m \ra \la k,a_n \ra = \lambda \la n+k,a_m \ra \la k,a_n \ra.
        \end{array}\right.
    \end{equation*}
    And so 
    \begin{equation*}
        \lambda \la n+k,a_m \ra\la n,a_k \ra \la m,a_n \ra = \lambda \la n,a_m \ra  (\la m,a_n \ra \la n,a_k \ra - \la m,a_k \ra \la k,a_n \ra ),
    \end{equation*}
    which leads to 
    \begin{equation} \label{eq : T3 condition colinearity}
        \la k,a_m \ra \la n,a_k \ra \la m,a_n \ra = - \la n,a_m \ra \la m,a_k \ra \la k,a_n \ra.
    \end{equation}
    Lets us prove that this cannot be verified in this case. By computation we obtain that 
    \begin{align*}
        \det\begin{pmatrix}
            a_m^1 &a_m^2 & a_m^3 \\
            a_n^1 & a_n^2 & a_n^3\\
            a_k^1 & a_k^2 & a_k^3
        \end{pmatrix}\begin{pmatrix}
            m^1 & n^1 &k^1 \\
            m^2 & n^2 &k^2 \\
            m^3 & n^3 &k^3 
        \end{pmatrix}&=\det\begin{pmatrix}
            0& \la n,a_m\ra & \la k,a_m\ra \\
            \la m,a_n \ra & 0 & \la k,a_n \ra \\
            \la m,a_k \ra & \la n,a_k \ra & 0
        \end{pmatrix} \\
        &=\la n,a_m\ra \la k,a_n \ra \la m,a_k \ra + \la k,a_m \ra \la m,a_n \ra \la n,a_k \ra.
    \end{align*}
    Equation \eqref{eq : T3 condition colinearity} is verified if and only if this determinant is equal to zero, which is impossible because $\mathrm{span}\left\{a_m,a_n,a_k\right\}=\mathrm{span}\left\{m,n,k\right\}=\R^3$. Then necessarily $p_{m,n,k}\we p_{n,k,m}\neq 0$. 
    
    In both cases we conclude thanks to \cref{generation of modes : from 2 non colinear modes}.
    \end{proof}

\end{document}
